\section{The adelic structure of a number field}

For a number field $K$, write $\mathcal O_K$ for its ring of integers.
A finite place $v_{\mathfrak p}$ is the normalized absolute value
$|a|_{\mathfrak p}=N(\mathfrak p)^{-\operatorname{ord}_{\mathfrak p}(a)}$
for $a\ne0$, and zero at zero, associated with a nonzero prime ideal
$\mathfrak p$ of $\mathcal O_K$.
An infinite place is the absolute value $a\mapsto|\sigma(a)|$ of a complex
embedding, with conjugate embeddings identified. Its multiplicity is one
for a real embedding and two for a nonreal conjugate pair.

We use the following external number-field results from
\cite{mathlib-number-fields}: $\mathcal O_K$ is a Dedekind domain, is finite
as a $\mathbb Z$-module, and has fraction field $K$; nonzero prime ideals
correspond bijectively to finite places; an element of $\mathcal O_K$ has
$\mathfrak p$-adic absolute value less than one precisely when it belongs
to $\mathfrak p$; there are finitely many infinite places; and two infinite
places have equivalent absolute values precisely when they are equal.
We also use the local denominator theorem: if
$\operatorname{ord}_{\mathfrak p}(y)\geq0$, then $y=m/d$ for
$m,d\in\mathcal O_K$ with $d\notin\mathfrak p$. The finite-support theorem
for finite places, lying over for integral extensions, and the number-field
product formula will be invoked explicitly below. These are inputs to the
construction, with the stated normalized absolute values.

\subsection{Places and their types}

\begin{definition}[Places]\label{def:places}
\lean{AbsoluteValue.Place,AbsoluteValue.trivialPlace}
For any field $F$, two real absolute values $v,w$ are equivalent if
$w=v^c$ for some $c>0$. A place of $F$ is an equivalence class of real
absolute values; write $\mathcal P_F$ for their set and $[v]$ for the class
of $v$. Write $v_\omega$ for a representative of $\omega\in\mathcal P_F$.
The trivial place $\omega_0$ is the class of the absolute value that is
zero at zero and one at every nonzero element.
\end{definition}

\begin{lemma}[Invariance of non-Archimedean type]\label{lem:place-type-invariance}
\uses{def:places}
\lean{NumberField.isNonarchimedean_congr}
Equivalent absolute values on any field are either both non-Archimedean
or both Archimedean.
\end{lemma}
\begin{proof}
Suppose $w=v^c$ with $c>0$ and $v(x+y)\leq\max(v(x),v(y))$.
The increasing function $t\mapsto t^c$ on $[0,\infty)$ gives
\[
 w(x+y)\leq\max(v(x)^c,v(y)^c)=\max(w(x),w(y)).
\]
Apply the same argument to $v=w^{1/c}$ for the converse.
\end{proof}

\begin{definition}[Types of places]\label{def:place-types}
\uses{def:places,lem:place-type-invariance}
\lean{AbsoluteValue.IsNonarchimedeanPlace,AbsoluteValue.IsArchimedeanPlace,AbsoluteValue.isNonarchimedeanPlace_or_isArchimedeanPlace}
A place is non-Archimedean if its representative satisfies the strong
triangle inequality, and Archimedean otherwise. This is independent of
the representative, and every place has exactly one of these two types.
\end{definition}
\begin{proof}
Representative independence follows from
Lemma~\ref{lem:place-type-invariance}. The two types are complementary
by their definitions.
\end{proof}

\begin{lemma}[The trivial place]\label{lem:trivial-place}
\uses{def:places,def:place-types}
\lean{AbsoluteValue.trivialPlace_out_eq,AbsoluteValue.isNonarchimedeanPlace_trivialPlace}
Every representative of $\omega_0$ is the trivial absolute value.
Equivalently, $v_\omega$ is trivial if and only if $\omega=\omega_0$.
The trivial place is non-Archimedean.
\end{lemma}
\begin{proof}
Every positive power of the trivial absolute value is itself; equivalence
therefore forces equality. Conversely a trivial representative has the
trivial equivalence class. To check the strong triangle inequality, if
both summands vanish then so does their sum. Otherwise the maximum of
their trivial absolute values is one, and the absolute value of their sum
is at most one.
\end{proof}

\begin{lemma}[A small element at a finite place]\label{lem:finite-place-small-element}
\lean{ArakelovGeometry.AdelicStructure.finitePlace_exists_lt_one}
For every finite place $v$ of $K$, some $a\in K^\times$ satisfies $v(a)<1$.
\end{lemma}
\begin{proof}
The nonzero prime ideal associated with $v$ contains a nonzero element
$a\in\mathcal O_K$. Its image in the fraction field is nonzero, and
membership in that ideal gives $v(a)<1$ by the finite-place membership
criterion.
\end{proof}

\begin{lemma}[Type of a finite place]\label{lem:finite-place-type}
\uses{def:places,def:place-types,lem:trivial-place,lem:place-type-invariance,lem:finite-place-small-element}
\lean{NumberField.finitePlace_isNonarchimedean,ArakelovGeometry.AdelicStructure.mk_finitePlace_ne_trivialPlace,ArakelovGeometry.AdelicStructure.isNonarchimedeanPlace_mk_finitePlace}
A normalized finite absolute value is non-Archimedean, and its class is
a nontrivial non-Archimedean place.
\end{lemma}
\begin{proof}
The valuation inequality
$\operatorname{ord}_{\mathfrak p}(a+b)\geq
\min(\operatorname{ord}_{\mathfrak p}(a),\operatorname{ord}_{\mathfrak p}(b))$
gives the strong triangle inequality after exponentiation by
$N(\mathfrak p)^{-1}<1$. Equivalent representatives have the same type.
If the class were trivial, its absolute value would be trivial by
Lemma~\ref{lem:trivial-place}, contradicting the nonzero small element
of Lemma~\ref{lem:finite-place-small-element}.
\end{proof}

\begin{lemma}[The value of two at an infinite place]\label{lem:infinite-place-two}
\lean{ArakelovGeometry.AdelicStructure.infinitePlace_apply_two}
For an infinite place $w$ of any field admitting a complex embedding,
$w(2)=2$.
\end{lemma}
\begin{proof}
Choose an embedding $\sigma$ representing $w$. A field embedding sends
$2$ to $2$, so $w(2)=|\sigma(2)|=2$.
\end{proof}

\begin{lemma}[Type of an infinite place]\label{lem:infinite-place-type}
\uses{def:place-types,lem:trivial-place,lem:place-type-invariance,lem:infinite-place-two}
\lean{NumberField.infinitePlace_not_isNonarchimedean,ArakelovGeometry.AdelicStructure.mk_infinitePlace_ne_trivialPlace,ArakelovGeometry.AdelicStructure.not_isNonarchimedeanPlace_mk_infinitePlace}
A normalized infinite absolute value on $K$ is Archimedean, and its class
is a nontrivial Archimedean place.
\end{lemma}
\begin{proof}
A non-Archimedean absolute value satisfies
$w(2)=w(1+1)\leq\max(w(1),w(1))=1$, whereas an infinite place has
$w(2)=2$. Thus it is Archimedean, as are its equivalent representatives.
It cannot be trivial because $2\ne0$ in $K$ and the trivial value of $2$
is one.
\end{proof}

\subsection{Classification of number-field places}

\begin{lemma}[Algebraic integers have bounded non-Archimedean absolute value]
\label{lem:integral-elements-bound}
A non-Archimedean absolute value $v$ on a field satisfies $v(a)\leq1$
for every element $a$ integral over $\mathbb Z$.
\end{lemma}
\begin{proof}
Repeated use of the strong triangle inequality gives $v(n)\leq1$ for
$n\in\mathbb Z$. Choose a monic equation
$a^d+\sum_{j<d}c_ja^j=0$, where $d\geq1$ and $c_j\in\mathbb Z$.
If $v(a)>1$, then
\[
 v(a)^d\leq\max_{j<d}v(c_j)v(a)^j
 \leq\max_{j<d}v(a)^j<v(a)^d,
\]
a contradiction.
\end{proof}

\begin{lemma}[Finite-place classification]\label{lem:number-field-nonarchimedean-classification}
\uses{def:places,lem:integral-elements-bound}
Every nontrivial non-Archimedean absolute value on $K$ is equivalent to
a normalized finite place.
\end{lemma}
\begin{proof}
Let $v$ be such an absolute value and put
$\mathfrak p=\{a\in\mathcal O_K:v(a)<1\}$. Every element of
$\mathcal O_K$ has value at most one by
Lemma~\ref{lem:integral-elements-bound}. The strong triangle inequality
shows that $\mathfrak p$ is closed under addition; multiplicativity and
the bound on $\mathcal O_K$ show that it is an ideal. It is proper because
$v(1)=1$. If $a,b\notin\mathfrak p$, then $v(a)=v(b)=1$, so
$ab\notin\mathfrak p$. Hence it is prime.

Nontriviality gives an $x\ne0$ with $v(x)<1$: invert an element whose
value is greater than one when necessary. Write $x=a/b$ with
$a,b\in\mathcal O_K$ and $b\ne0$. Then $a\ne0$ and
$v(a)=v(x)v(b)\leq v(x)<1$, so $\mathfrak p\ne0$.

Let $u=v_{\mathfrak p}$. If $u(y)<1$, the local denominator theorem
gives $y=m/d$ with $d\notin\mathfrak p$ and $m,d\in\mathcal O_K$.
Since $\operatorname{ord}_{\mathfrak p}(d)=0$ and
$\operatorname{ord}_{\mathfrak p}(y)>0$, one has $m\in\mathfrak p$.
Consequently $v(d)=1$, $v(m)<1$, and $v(y)<1$. The absolute-value
comparison theorem says that for a nontrivial absolute value $u$, the
implication $u(y)<1\Rightarrow v(y)<1$ for all $y$ makes $u$ and $v$
equivalent. Applying this external comparison theorem proves the claim;
$u$ is nontrivial because it is less than one on the nonzero element
$a\in\mathfrak p$ constructed above.
\end{proof}

\begin{lemma}[Archimedean classification]\label{lem:number-field-archimedean-classification}
\uses{def:places,lem:completion-nonarchimedean,thm:complete-archimedean}
Every Archimedean absolute value on $K$ is equivalent to a normalized
infinite place.
\end{lemma}
\begin{proof}
The absolute value on the completion of $K$ remains Archimedean: a strong
triangle inequality on the completion would restrict to one on $K$.
Theorem~\ref{thm:complete-archimedean} identifies that completion with
$\mathbb R$ or $\mathbb C$, with its absolute value equal to the $s$th
power of the usual modulus for some $s>0$. Compose the inclusion of $K$
into its completion with this identification and, in the real case, with
$\mathbb R\hookrightarrow\mathbb C$. The resulting embedding $\sigma$
satisfies $v(a)=|\sigma(a)|^s$ for every $a\in K$. Its infinite place is
therefore equivalent to $v$.
\end{proof}

\begin{theorem}[Ostrowski classification for number fields]\label{thm:number-field-ostrowski}
\uses{lem:number-field-nonarchimedean-classification,lem:number-field-archimedean-classification}
\lean{NumberField.ostrowski}
Every nontrivial absolute value on $K$ is equivalent to a normalized
finite place or a normalized infinite place.
\end{theorem}
\begin{proof}
The nonarchimedean case follows from
Lemma~\ref{lem:number-field-nonarchimedean-classification}, and
Lemma~\ref{lem:number-field-archimedean-classification} proves the
archimedean case. These cases exhaust all absolute values.
\end{proof}

\begin{proposition}[Classification of places]\label{prop:number-field-place-classification}
\uses{thm:number-field-ostrowski,lem:place-type-invariance,lem:trivial-place,lem:finite-place-type,lem:infinite-place-type}
\lean{NumberField.isNonarchimedeanPlace_iff,NumberField.isArchimedeanPlace_iff,ArakelovGeometry.AdelicStructure.place_trichotomy}
A nontrivial place $\omega$ is non-Archimedean if and only if
$v_\omega$ is equivalent to a finite place. A place $\omega$ is
Archimedean if and only if $v_\omega$ is equivalent to an infinite place.
Every place is therefore trivial, the class of a finite place, or the
class of an infinite place, and these three cases are disjoint.
\end{proposition}
\begin{proof}
A representative of a nontrivial place is nontrivial, since equality
with the trivial absolute value would give the trivial class.
Theorem~\ref{thm:number-field-ostrowski} applies to such a representative.
Equivalence preserves non-Archimedean type, while finite and infinite
places have opposite types, so the two alternatives give precisely the
two asserted characterizations. An Archimedean representative cannot be
trivial by Lemma~\ref{lem:trivial-place}. Equality of equivalence classes
then gives the trichotomy. The finite and infinite classes are nontrivial
and have opposite types, proving disjointness.
\end{proof}

\begin{lemma}[Finite places give distinct classes]\label{lem:finite-place-injectivity}
\uses{def:places}
\lean{NumberField.finitePlace_injective_place}
The map $v\mapsto[v]$ from finite places to $\mathcal P_K$ is injective.
\end{lemma}
\begin{proof}
Equivalent absolute values have the same elements of value less than one,
since a positive power preserves that inequality. Applied to
$a\in\mathcal O_K$, the finite-place membership criterion shows that
two equivalent finite places have the same prime ideal. The bijection
between nonzero prime ideals and normalized finite places then makes
them equal.
\end{proof}

\subsection{Countability and finite support}

\begin{lemma}[Countability of places]\label{lem:number-field-places-countable}
\uses{prop:number-field-place-classification}
\lean{NumberField.instCountablePlace}
The set $\mathcal P_K$ is countable.
\end{lemma}
\begin{proof}
The ring $\mathcal O_K$ is a finitely generated $\mathbb Z$-module,
so it is countable. It is Noetherian, hence every ideal is generated by
a finite subset. There are countably many finite subsets of a countable
set, so there are countably many ideals, nonzero prime ideals, and finite
places. There are finitely many infinite places. Encode a place by a
distinguished symbol in the trivial case, a selected finite place in the
nontrivial non-Archimedean case, and a selected infinite place in the
Archimedean case. This maps into a countable disjoint union. It is
injective: equality of the codes forces the same case and the same
selected absolute value, hence the same equivalence class. Thus
$\mathcal P_K$ is countable.
\end{proof}

\begin{lemma}[Infinitely many nonzero prime ideals]\label{lem:number-field-primes-infinite}
\lean{NumberField.infinite_heightOneSpectrum}
The ring $\mathcal O_K$ has infinitely many nonzero prime ideals.
\end{lemma}
\begin{proof}
For each rational prime $p$, lying over for the integral extension
$\mathbb Z\subseteq\mathcal O_K$ gives a prime ideal $\mathfrak p_p$
whose contraction is $(p)$. It is nonzero, since the zero ideal contracts
to zero under the inclusion and $(p)\ne0$. If
$\mathfrak p_p=\mathfrak p_q$, contraction gives $(p)=(q)$ in
$\mathbb Z$, so $p=\pm q$. Rational primes are positive, hence $p=q$.
Thus the infinitely many rational primes inject into the nonzero prime
ideals of $\mathcal O_K$.
\end{proof}

\begin{corollary}[Infinitude of places]\label{cor:number-field-places-infinite}
\uses{lem:number-field-primes-infinite,lem:finite-place-injectivity}
\lean{NumberField.instInfinitePlace}
The set $\mathcal P_K$ is infinite.
\end{corollary}
\begin{proof}
Nonzero prime ideals give infinitely many finite places, and distinct
finite places give distinct elements of $\mathcal P_K$ by
Lemma~\ref{lem:finite-place-injectivity}.
\end{proof}

\begin{lemma}[Finiteness of the Archimedean locus]\label{lem:number-field-archimedean-finite}
\uses{prop:number-field-place-classification}
\lean{NumberField.finite_isArchimedeanPlace}
The Archimedean places form a finite subset of $\mathcal P_K$.
\end{lemma}
\begin{proof}
Select an infinite place representing each Archimedean place by the
classification. If two selected infinite places coincide, their classes
coincide, so the selection is injective. Its target is the finite set of
infinite places, proving the assertion.
\end{proof}

\begin{lemma}[Finitely many non-Archimedean contributions]\label{lem:number-field-finite-contributions}
\uses{prop:number-field-place-classification,lem:trivial-place}
\lean{NumberField.finite_isNonarchimedean_ne_one}
For $a\in K^\times$, only finitely many non-Archimedean places $\omega$
satisfy $v_\omega(a)\ne1$.
\end{lemma}
\begin{proof}
The trivial place has value one at $a$, so every place in question is
represented by a finite place $v$. A positive power takes the value one
exactly at one; hence $v(a)\ne1$. The external finite-support theorem
for normalized finite places says that $\{v:v(a)\ne1\}$ is finite.
Every place in the statement is the class of an element of this finite
set. Selecting such a representative gives an injection into that set,
because identical selected representatives have identical classes.
There are therefore only finitely many places in question.
\end{proof}

\begin{lemma}[Finite support of logarithmic evaluations]\label{lem:number-field-finite-support}
\uses{def:place-types,lem:number-field-archimedean-finite,lem:number-field-finite-contributions}
\lean{NumberField.finite_out_ne_one}
For $a\in K^\times$, the set $\{\omega:v_\omega(a)\ne1\}$ is finite.
\end{lemma}
\begin{proof}
It is contained in the union of the finite Archimedean locus and the
finite set of non-Archimedean places contributing a value different from
one. Both sets are finite by the preceding two lemmas.
\end{proof}

\subsection{Normalized representatives and the adelic measure}

\begin{definition}[Normalized representatives and weights]\label{def:number-field-place-data}
\uses{prop:number-field-place-classification}
\lean{ArakelovGeometry.AdelicStructure.placeRep,ArakelovGeometry.AdelicStructure.placeWeight}
For $\omega\in\mathcal P_K$ choose the normalized representative
$r_\omega$ as follows: at $\omega_0$ take the trivial absolute value;
at a nontrivial non-Archimedean place take a finite place in its class;
at an Archimedean place take an infinite place in its class. Put
\[
 c(\omega)=
 \begin{cases}
 0,&\omega=\omega_0,\\
 1,&\omega\text{ is nontrivial and non-Archimedean},\\
 m(w),&\omega=[w]\text{ for the chosen infinite place }w.
 \end{cases}
\]
Here $m(w)$ is its multiplicity, one or two.
\end{definition}

\begin{lemma}[The chosen representative belongs to its place]\label{lem:number-field-representative-equivalence}
\uses{def:number-field-place-data,lem:trivial-place}
\lean{ArakelovGeometry.AdelicStructure.placeRep_equiv}
For every place $\omega$, the representative $v_\omega$ is equivalent
to $r_\omega$.
\end{lemma}
\begin{proof}
At the trivial place the representatives are equal. In each of the two
other cases, equivalence is the defining property of the finite or infinite
place selected from the classification.
\end{proof}

\begin{lemma}[Recovery of finite-place data]\label{lem:number-field-finite-normalization}
\uses{def:number-field-place-data,lem:finite-place-type,lem:finite-place-injectivity}
\lean{ArakelovGeometry.AdelicStructure.finitePlace_choose_eq,ArakelovGeometry.AdelicStructure.placeRep_mk_finitePlace,ArakelovGeometry.AdelicStructure.placeWeight_mk_finitePlace}
For a finite place $v$, any finite place selected to represent $[v]$ is
$v$ itself. Consequently $r_{[v]}=v$ and $c([v])=1$.
\end{lemma}
\begin{proof}
The class $[v]$ is nontrivial and non-Archimedean. Thus the finite-place
case of Definition~\ref{def:number-field-place-data} applies. Its chosen
place has the same class as $v$, so injectivity in
Lemma~\ref{lem:finite-place-injectivity} makes it equal to $v$.
The representative and weight formulas follow from their definitions.
\end{proof}

\begin{lemma}[Recovery of infinite-place data]\label{lem:number-field-infinite-normalization}
\uses{def:number-field-place-data,lem:infinite-place-type}
\lean{ArakelovGeometry.AdelicStructure.infinitePlace_choose_eq,ArakelovGeometry.AdelicStructure.placeRep_mk_infinitePlace,ArakelovGeometry.AdelicStructure.placeWeight_mk_infinitePlace}
For an infinite place $w$, any infinite place selected to represent $[w]$
is $w$ itself. Consequently $r_{[w]}=w$ and $c([w])=m(w)$.
\end{lemma}
\begin{proof}
The place $[w]$ is Archimedean and nontrivial. Thus the infinite-place
case of Definition~\ref{def:number-field-place-data} applies. The selected
place is equivalent to $w$; the external uniqueness theorem for normalized
infinite places makes them equal. Substitution gives both asserted formulas.
\end{proof}

\begin{lemma}[Finite weights]\label{lem:number-field-weights}
\uses{def:number-field-place-data}
\lean{ArakelovGeometry.AdelicStructure.placeWeight_ne_top,ArakelovGeometry.AdelicStructure.placeWeight_trivialPlace}
Every $c(\omega)$ is finite, and $c(\omega_0)=0$.
\end{lemma}
\begin{proof}
The possible weights are zero, one, and the finite multiplicity of an
infinite place. The trivial-place value is the first case of the definition.
\end{proof}

\begin{definition}[The number-field adelic structure]\label{def:number-field-adelic-structure}
\uses{def:adelic-structure,def:singleton-measure,def:number-field-place-data,lem:number-field-representative-equivalence,lem:number-field-weights,lem:number-field-places-countable,lem:number-field-finite-support}
\lean{ArakelovGeometry.AdelicStructure.numberFieldAdelicStructure,ArakelovGeometry.AdelicCurve.numberFieldAdelicCurve}
Equip $\mathcal P_K$ with its discrete sigma algebra and the measure
\[
 \nu_K=\sum_{\omega\in\mathcal P_K}c(\omega)\delta_\omega.
\]
Together with $\omega\mapsto r_\omega$, this is an adelic structure on
$K$, and hence an adelic curve with underlying field $K$.
\end{definition}
\begin{proof}
The sum of nonnegative multiples of Dirac measures is a measure. Every
function on the discrete measurable space is measurable. Fix $a\ne0$.
If $v_\omega(a)=1$, equivalence gives $r_\omega(a)=1$, so
$\log r_\omega(a)=0$. Thus the logarithmic evaluation vanishes outside
the finite set in Lemma~\ref{lem:number-field-finite-support}.
The Dirac-sum integrability criterion on a countable discrete space says
that a real function $f$ is integrable precisely when
$\sum_\omega c(\omega)|f(\omega)|<\infty$, provided the weights are
finite. Here the summands are finite and supported on a finite set, so
that criterion proves integrability. These are the defining requirements
of an adelic structure.
\end{proof}

\begin{theorem}[The logarithmic product formula]\label{thm:number-field-product-formula}
\uses{def:number-field-adelic-structure,prop:number-field-place-classification,lem:finite-place-type,lem:infinite-place-type,lem:finite-place-injectivity,lem:number-field-finite-normalization,lem:number-field-infinite-normalization,lem:number-field-weights}
\lean{ArakelovGeometry.AdelicCurve.numberField_product_formula}
For every $a\in K^\times$,
\[
 \int_{\mathcal P_K}\log r_\omega(a)\,d\nu_K(\omega)=0.
\]
\end{theorem}
\begin{proof}
The integral formula for a weighted sum of Dirac measures gives
\[
 \int\log r_\omega(a)\,d\nu_K
 =\sum_\omega c(\omega)\log r_\omega(a).
\]
This sum has finite support, as proved in
Definition~\ref{def:number-field-adelic-structure}. The trivial place
contributes zero. The finite and infinite classes are disjoint; finite
places inject into their classes by Lemma~\ref{lem:finite-place-injectivity},
and infinite places do so by the uniqueness theorem for their normalized
absolute values. The classification therefore lets us reindex the sum as
\[
 \sum_{w\text{ infinite}}m(w)\log w(a)
 +\sum_{v\text{ finite}}\log v(a).
\]
The second sum has finitely many nonzero terms by the finite-support
theorem for finite places. Every value in these sums is positive because
$a\ne0$. The rules for the logarithm of a finite product therefore turn
the expression into
\[
 \log\left(
   \prod_{w\text{ infinite}}w(a)^{m(w)}
   \prod_{v\text{ finite}}v(a)\right).
\]
The external number-field product formula in
\cite{mathlib-number-fields}, with exactly these normalizations, makes
the product one. Its logarithm is zero, as required.
\end{proof}

\begin{corollary}[Properness of the number-field adelic curve]\label{cor:number-field-proper}
\uses{def:adelic-structure,def:number-field-adelic-structure,thm:number-field-product-formula}
The adelic curve associated with $K$ is proper.
\end{corollary}
\begin{proof}
Theorem~\ref{thm:number-field-product-formula} gives zero for the integral
of every nonzero logarithmic evaluation. This is precisely properness.
\end{proof}
